Adaptive Dynamic Programming (ADP) approximates dynamic programming when the exact value function is unavailable or too expensive to compute. In the actor-critic form, the critic learns a value approximation $\hat V_\theta(x)$ and the actor updates a policy using the critic's gradient. This is closely related to reinforcement learning, but in continuous-time optimal control the update is usually motivated by the HJB equation and policy iteration~\cite{bertsekas1996neuro,sutton2018reinforcement}.

For a fixed policy $u=\pi(x)$, the critic is trained to satisfy the policy-evaluation equation
\begin{equation}
\begin{aligned}
0 ={}& c(x,\pi(x)) + \nabla V^\pi(x)^\top f(x,\pi(x))\\
&+ \frac{1}{2}\operatorname{Tr}\left(\sigma^\top \nabla^2 V^\pi(x)\sigma\right).
\end{aligned}
\end{equation}
The actor then improves the policy by greedily minimizing the Hamiltonian,
\begin{equation}
\pi_{\mathrm{new}}(x)=\arg\min_u
\left[c(x,u)+\nabla \hat V_\theta(x)^\top f(x,u)\right].
\end{equation}
For the LQR case study, this greedy update has the closed form
\begin{equation}
u^\star(x)=-\frac{1}{2}R^{-1}B^\top \nabla \hat V_\theta(x).
\end{equation}

The Julia implementation uses the quadratic critic
\begin{equation}
\hat V_\theta(x)=x^\top P_\theta x+\beta_\theta,
\end{equation}
then alternates between policy evaluation for the current linear controller and policy improvement using the value gradient. This is the most structure-aware method in the comparison: it uses the simulator, the cost, and the known control-affine dynamics without requiring a full PDE grid. Its main limitation is that the critic class must be expressive enough; if the true value is far from quadratic, the same actor-critic loop would need a richer function approximator and more careful stabilization.

\paragraph{Convergence requirements.}
Policy iteration converges monotonically to $V^\star$ when the initial policy
is \emph{admissible}---stabilizing with finite cost---a hard prerequisite that
must be constructed by the user~\cite{abu2005nearly}.
With function approximation, convergence is to a neighborhood of $V^\star$
whose size depends on the expressiveness of the basis class.
The online simultaneous actor-critic of~\cite{vamvoudakis2010online} provides
a Lyapunov proof of \emph{uniform ultimate boundedness} of the weight errors,
not asymptotic convergence to $V^\star$.
For the stochastic case, the policy-evaluation equation~\eqref{eqn:stochastic_dynamics}
requires the Ito correction term $\frac{1}{2}\operatorname{Tr}(\sigma^\top\nabla^2 V^\pi\sigma)$;
basis functions must be chosen so that this term is representable in the critic class.
